\scanpage{49}
% archive sheet number: 48
est majoré par $(n-1)d$, (au lieu de $(n-1)\unsure{n^2|n}$) où $d=\dim G$. \deleted{\unsure{Les|Ces} \unsure{$g$}} D'ailleurs, si $v_1,v_2,\ldots v_n$ sont donnés, \deleted{\unsure{alors|les}} \deleted{\unsure{permettent|peuvent}} définissant $g_2,\ldots,g_n\in G(\widehat{\mathbb Z})$
alors prenant $\sigma_2\ldots\sigma_n\in\operatorname{Gal}(k/k_0)$ (où $k_0$ est un corps de définition \inserted{de type fini} de $M$) et $v'_1,v'_2,\ldots v'_n$ avec $v'_i=\sigma_i v_i\quad 1\leq i\leq n$, on trouve les valeurs \deleted{$g$} $\dot{\sigma}_2g_2,\ldots\dot{\sigma}_ng_n$, où $\dot{\sigma}_i$ est l'image dans $G(\deleted{\unsure{\widehat{\mathbb Q}}}\widehat{\mathbb Z})$. \deleted{\unsure{Les}}

\margin{N.B. On \unsure{prendra} le $G$ ``\uline{\unsure{géom|vrai|gros|fine}}'' donc \uline{connexe}}
% slanted
Admettant la \uuline{conjecture de Tate renforcée}, $\operatorname{Gal}(k/k_0)\unsure{\mathrel{\text{---}}|\longrightarrow}G(\widehat{\mathbb Z})$ \uline{a une image qui est d'indice fini.}
Il s'ensuit \deleted{que} quitte à remplacer les $v_i$ par des $\sigma_i v_i$ convenables, on peut \deleted{\unsure{trouver}} \inserted{obtenir} $g_2\ldots g_n$ gén. indépendants pour \uline{tt} $l$, si on admet ceci :

\[
\left]\begin{aligned}
&\text{Soit }G\ \text{gpe alg. affine connexe sur }\mathbb Q\\
&U\ \text{sous-gpe ouvert cpt de }G(\widehat{\mathbb Q})\\
&P\inserted{=U.a}\ \text{une classe en }G(\widehat{\mathbb Q})\ \text{mod }U\\
&\text{Alors il existe dans }P\ \text{un élément }\unsure{\xi|\zeta}\ \text{dont toutes les}\\
&\hspace{3cm}\text{composantes sont génériques rel./}\mathbb Q
\end{aligned}\right.
\]
